% !PRINTED 737
\weberpage{737}

Hiernach ergiebt sich der Regulator $R$ des Systems (7) aus
der Determinante:

\[
\left|
\begin{array}{cccccccc}
l'_{1,1}, & l'_{2,1}, & \ldots, & l'_{\nu'-1,1}, & L_{0,1}, & L_{1,1}, & \ldots, & L_{\nu'-1,1} \\
\cdots & \cdots & \cdots & \cdots & \cdots & \cdots & \cdots & \cdots \\
l'_{1,\nu'-1}, & l'_{2,\nu'-1}, & \ldots, & l'_{\nu'-1,\nu'-1}, & L_{0,\nu'-1}, & L_{1,\nu'-1}, & \ldots, & L_{\nu'-1,\nu'-1} \\
l'_{1,0}, & l'_{2,0}, & \ldots, & l'_{\nu'-1,0}, & -L_{0,0}, & -L_{1,0}, & \ldots, & -L_{\nu'-1,0} \\
l'_{1,1}, & l'_{2,1}, & \ldots, & l'_{\nu'-1,1}, & -L_{0,1}, & -L_{1,1}, & \ldots, & -L_{\nu'-1,1} \\
\cdots & \cdots & \cdots & \cdots & \cdots & \cdots & \cdots & \cdots \\
l'_{1,\nu'-1}, & l'_{2,\nu'-1}, & \ldots, & l'_{\nu'-1,\nu'-1}, & -L_{0,\nu'-1}, & -L_{1,\nu'-1}, & \ldots, & -L_{\nu'-1,\nu'-1}
\end{array}
\right|.
\]

Diese Determinante hat $\nu-1$ Zeilen; wir addiren die erste
zur $(\nu'+1)^{\mathrm{ten}}$, die zweite zur $(\nu'+2)^{\mathrm{ten}}$ u.\,s.\,f., die $(\nu'-1)^{\mathrm{te}}$ zur
letzten. Dadurch heben sich in den letzten $(\nu'-1)$ Zeilen die
$L$ heraus und die $l$ bekommen den Factor 2 [nach (10), (11)],
und hieraus ergiebt sich nach einem Determinantensatze (Bd.\,I,
\S.\,23, V.):
\weq{12}{R = 2^{\nu'-1} E' T_0.}

Da hiernach $R$ nicht verschwindet und daher in (7) ein
System von $\nu-1$ unabhängigen Einheiten des Körpers $H_m$ vorliegt,
so können wir [nach \S.\,186, (5)] die rationalen (ganzen
oder gebrochenen) Zahlen $m_{i,k}$, $M_{i,k}$ aus den Gleichungen
\weq{13}{
\begin{aligned}
2\, l_{i,k} & = m_{1,i}\, l'_{1,k} + \cdots + m_{\nu'-1,i}\, l'_{\nu'-1,k} \\
& \quad + M_{0,i}\, L_{0,k} + M_{1,i}\, L_{1,k} + \cdots + M_{\nu'-1,i}\, L_{\nu'-1,k}
\end{aligned}
}

bestimmen. Es lässt sich aber leicht durch die Relationen (10), (11)
nachweisen, dass die Zahlen $m_{i,k}$, $M_{i,k}$ hier \gesperrt{ganze} sein müssen.
Denn es ist
\[
\begin{aligned}
l_{i,k} + l_{i,k+\nu'} & = \log \lvert \varepsilon_i(r_k)\, \varepsilon_i(-r_k) \rvert = \\
& \quad m_{1,i}\, l'_{1,k} + \cdots + m_{\nu'-1,i}\, l'_{\nu'-1,k} \\
l_{i,k} - l_{i,k+\nu'} & = \log \lvert \varepsilon_i(r_k)\, \varepsilon_i(-r_k)^{-1} \rvert = \\
& \quad M_{0,i}\, L_{0,k} + M_{1,i}\, L_{1,k} + \cdots + M_{\nu'-1,i}\, L_{\nu'-1,k}.
\end{aligned}
\]

Die Einheiten $\varepsilon_i(r)\, \varepsilon_i(-r)$ gehören aber dem Körper $H_\mu$
an, und da die $\varepsilon'_i$ nach Voraussetzung ein Fundamentalsystem
dieses Körpers sind, so müssen die $m_{i,k}$ nach \S.\,187 ganze rationale
 Zahlen sein.

Die Einheiten $\varepsilon_i(r)\, \varepsilon_i(-r)^{-1}$ sind Normaleinheiten des Körpers
 $H_m$, und weil die $\mathcal{E}_i$ ein Fundamentalsystem von Normaleinheiten
 sein sollten, so sind nach \S.\,200 die $M_{i,k}$ ganze Zahlen.
